\subsection{问题一结果}

程序构造了一个确由三次示向楔形交得到的等边三角形可行域，顶点为
$(0,0)$、$(40,0)$ 和 $(20,20\sqrt3)$。由式\eqref{eq:diameter}得
$D=40\unit{m}$。若以底边为直径作半径 $20\unit{m}$ 的圆，第三个顶点位于圆外；该区域的最小包围圆圆心为
$(20,20/\sqrt3)$，半径为
\begin{equation}
r_*=\frac{40}{\sqrt3}=23.094\unit{m}>\frac D2.
\end{equation}
因此问题一的覆盖判断为否。图\ref{fig:q1}同时给出直径圆与最小包围圆。
\begin{figure}[htbp]
  \centering
  \includegraphics[width=0.68\textwidth]{q1-region-counterexample.pdf}
  \caption{直径圆不能覆盖定位区域的可行示向反例}
  \label{fig:q1}
\end{figure}

\subsection{问题二结果}

将 $R=1500\unit{m}$、$r_0=1000\unit{m}$、
$\varepsilon=1^\circ$ 代入式\eqref{eq:second-point}，在以首次示向中央线为
$x$ 轴的局部坐标中得到
\begin{equation}
S_2^\pm-S_1=(761.429,\ \pm648.248)\unit{m}.
\end{equation}
两个点到 $S_1,E_-,E_+$ 的最大距离均为 $1000\unit{m}$，故属于候选域
$\mathcal C$，移动距离为 $1000\unit{m}$，对应 $200\unit{s}$。在示向角取
$[-1^\circ,1^\circ]$ 的 321 个网格、距离取 $(5,1500]$ 的 401 个网格上，最小交会角为 $39.592^\circ$；中央线上点 $(750,0)$ 虽能保证接收，但最小交会角为 $0^\circ$。纯垂向点 $(0,1000)$ 到可能源的最大距离为
$1817.239\unit{m}$，无法保证第二次收到信号。候选域和推荐点见图\ref{fig:q2}。

\begin{figure}[htbp]
  \centering
  \includegraphics[width=0.72\textwidth]{q2-second-detector-region.pdf}
  \caption{保证第二次接收的三圆盘候选域及对称推荐点}
  \label{fig:q2}
\end{figure}

\subsection{问题三、四演练结果}

截至本文稿冻结时，问题三、四各完成三轮官方演练。表\ref{tab:practice}中的“程序运行时间”为现实世界程序耗时，其余时间指标均为虚拟时间。六轮共清除 $82/82$ 个源，清除比例均为 1；全部请求被接受，没有清除失败。问题四三轮共包含 18 个定向源。

\begin{table}[htbp]
  \centering
  \caption{问题三、四官方演练结果（非正式成绩）}
  \label{tab:practice}
  \resizebox{\textwidth}{!}{%
  \begin{tabular}{clrrrrr}
    \toprule
    问题 & 测试案例编码 & 源数 & 清除数 & 清除比例 & 平均定位清除时间/s & 程序运行时间/s \\
    \midrule
    3 & AFYH-6GTT-GWWB-ZDPW & 11 & 11 & 1.000 & 1539.048 & 6.210 \\
    3 & FYKN-VY2Y-C99F-6BZN & 15 & 15 & 1.000 & 1113.106 & 4.368 \\
    3 & 8BS5-GQWW-6BTJ-8Z73 & 15 & 15 & 1.000 & 1085.527 & 4.432 \\
    4 & 77QS-523V-UX4E-UCGC & 14 & 14 & 1.000 & 1137.709 & 5.302 \\
    4 & UC47-56D3-EHEC-RTK3 & 16 & 16 & 1.000 & 1099.899 & 5.406 \\
    4 & F6QC-7NFR-B6ES-WJMF & 11 & 11 & 1.000 & 1503.504 & 6.645 \\
    \bottomrule
  \end{tabular}
  }
\end{table}

图\ref{fig:practice}给出逐轮清除比例及虚拟时间构成。移动时间在六轮中均占主要部分，说明在保持覆盖证明的前提下，优化路线续接比减少定位计算更能改善总时间。

\begin{figure}[htbp]
  \centering
  \includegraphics[width=0.98\textwidth]{q34-practice-summary.pdf}
  \caption{六轮官方演练的清除比例与虚拟时间构成}
  \label{fig:practice}
\end{figure}

把式\eqref{eq:time}按动作拆分可进一步判断主要瓶颈。表\ref{tab:practice-time}中的移动、检测、切频与清除四项之和均与模拟器虚拟总时间相符，最大差异来自小数舍入。六轮移动时间占比为 $73.66\%$--$83.03\%$；相比之下，成功清除本身只占 $55$--$80\unit{s}$。因此继续减少清除计算耗时几乎不会改变总成绩，路线长度和无效扫描才是主要优化对象。

\begin{table}[htbp]
  \centering
  \caption{六轮官方演练的虚拟时间分解（单位：秒）}
  \label{tab:practice-time}
  \resizebox{\textwidth}{!}{%
  \begin{tabular}{crrrrrr}
    \toprule
    轮次 & 移动 & 检测 & 切频 & 清除 & 虚拟总时间 & 移动占比/\% \\
    \midrule
    问题3--第1轮 & 12763.531 & 3520 & 591 & 55 & 16929.531 & 75.39 \\
    问题3--第2轮 & 13762.589 & 2490 & 369 & 75 & 16696.589 & 82.43 \\
    问题3--第3轮 & 13336.902 & 2495 & 376 & 75 & 16282.902 & 81.91 \\
    问题4--第1轮 & 12578.931 & 2835 & 444 & 70 & 15927.931 & 78.97 \\
    问题4--第2轮 & 14611.383 & 2535 & 372 & 80 & 17598.383 & 83.03 \\
    问题4--第3轮 & 12181.539 & 3680 & 622 & 55 & 16538.539 & 73.66 \\
    \bottomrule
  \end{tabular}}
\end{table}

六轮虚拟总时间均值为 $16662.312\unit{s}$，总体标准差为
$523.991\unit{s}$，最小值与最大值分别为 $15927.931\unit{s}$、
$17598.383\unit{s}$。不过各轮源数、定向比例和策略版本不同，这组离散程度只用于描述演练运行范围，不能当作同分布重复试验的置信区间。现实程序运行时间均不超过 $6.645\unit{s}$，表明几何运算和状态管理没有成为现实时间瓶颈。

同种子 100 局合成对照显示，从固定蛇形续扫改为“定位后前往最近未完成覆盖点”，问题三、四平均虚拟时间分别下降 $12.70\%$ 和 $9.99\%$，平均移动距离分别下降 $10.30\%$ 和 $8.56\%$，且仍为 100/100 全部清除。在 v3 中，达到题面上限 16 个源后立即结束；16 源子集相对 v2 的平均虚拟时间在问题三、四分别下降 $32.59\%$ 和 $19.60\%$。

\subsection{正式测试结果表}

正式测试尚未启动。为避免把演练数据冒充正式成绩，表\ref{tab:formal3}和表\ref{tab:formal4}只保留题目要求的结构；每轮完成后将从模拟器结果和原始日志自动填入。

\begin{table}[htbp]
  \centering
  \caption{问题三正式测试结果}
  \label{tab:formal3}
  \begin{tabular}{lccc}
    \toprule
    测试案例编码 & 清除干扰源个数 & 平均定位清除时间/s & 程序运行时间/s \\
    \midrule
    尚未执行（第1轮） & -- & -- & -- \\
    尚未执行（第2轮） & -- & -- & -- \\
    尚未执行（第3轮） & -- & -- & -- \\
    \bottomrule
  \end{tabular}
\end{table}

\begin{table}[htbp]
  \centering
  \caption{问题四正式测试结果}
  \label{tab:formal4}
  \begin{tabular}{lccc}
    \toprule
    测试案例编码 & 清除干扰源个数 & 平均定位清除时间/s & 程序运行时间/s \\
    \midrule
    尚未执行（第1轮） & -- & -- & -- \\
    尚未执行（第2轮） & -- & -- & -- \\
    尚未执行（第3轮） & -- & -- & -- \\
    \bottomrule
  \end{tabular}
\end{table}
